\chapter{Cascade Control Implementation}

\section{Why Cascade Control}

\subsection{Limitation of Single-Loop Control on $T_{out2}$}

A direct single-loop PID controller regulating $T_{out2,HE1}$ by manipulating $F_{1,SP}$
faces the following challenges:

\begin{enumerate}[leftmargin=*, itemsep=4pt]
  \item \textbf{Slow combined dynamics.}
    The path $F_{1,SP} \rightarrow T_{out2}$ passes through the furnace (time constant
    $\tau_1$) and then the heat exchanger ($\tau_2$). The total apparent time constant
    $\tau_{\text{total}} \approx \tau_1 + \tau_2$ forces a long integral time, producing
    sluggish setpoint tracking.

  \item \textbf{Late disturbance detection.}
    A change in furnace temperature ($T_{hin}$ disturbance) first affects $T_{in1}$ and
    only later, after traversing the HE dynamics, affects $T_{out2}$. The single-loop
    controller cannot correct the disturbance until it has propagated through both stages.

  \item \textbf{Stability--performance trade-off.}
    The phase delay of the combined second-order-like dynamics limits the achievable
    closed-loop bandwidth. Aggressive tuning leads to oscillations; conservative tuning
    leaves disturbances largely uncorrected.
\end{enumerate}

\subsection{Cascade Principle}

Cascade control (Figure~\ref{fig:cascade_block}) addresses these limitations by adding an
intermediate measurement loop:

\begin{itemize}[leftmargin=*, itemsep=3pt]
  \item \textbf{Inner loop (fast):} Regulates $T_{in1,HE1}$ by manipulating $F_{1,SP}$.
    Because $T_{in1}$ responds to $F_{1,SP}$ much faster than $T_{out2}$ does, the inner
    loop can be tuned aggressively.
  \item \textbf{Outer loop (slow):} Sets the $T_{in1}$ setpoint based on the $T_{out2}$
    error. Operates at a bandwidth commensurate with the HE dynamics ($\tau_2$).
\end{itemize}

\textbf{Key requirement:} The inner loop bandwidth must be 3--5$\times$ higher than the
outer loop bandwidth. For this installation, $T_{in1}$ responds to $F_{1,SP}$ faster than
$T_{out2}$ responds to $T_{in1}$ (the furnace is a direct thermal source; the HE has
distributed thermal mass), so the condition is satisfied.

\begin{figure}[H]
  \centering
  \begin{tikzpicture}[
    block/.style  = {rectangle, draw=darkblue, thick, fill=blue!8,
                     minimum width=2.4cm, minimum height=1.0cm,
                     font=\small, align=center},
    sum/.style    = {circle, draw=darkblue, thick, fill=white,
                     minimum size=0.65cm, inner sep=0pt, font=\footnotesize},
    signal/.style = {-Stealth, thick, draw=darkblue},
    label/.style  = {font=\scriptsize, above},
  ]
    % Nodes
    \node[sum] (s1) at (0,0)    {$\Sigma$};
    \node[block] (Co) at (2.0,0)  {$C_{\text{outer}}$\\$K_{p,ext}$};
    \node[sum] (s2) at (4.2,0)  {$\Sigma$};
    \node[block] (Ci) at (6.2,0)  {$C_{\text{inner}}$\\PI+P};
    \node[block] (G1) at (8.8,0)  {$G_1(s)$\\$F_{1,SP} \to T_{in1}$};
    \node[block] (G2) at (11.4,0) {$G_2(s)$\\$T_{in1} \to T_{out2}$};

    % Setpoint input
    \draw[signal] (-1.5,0) -- (s1) node[label, pos=0]{$T_{out2,SP}$};
    \node[font=\scriptsize] at (s1.north east) {$+$};
    \node[font=\scriptsize] at (s1.south east) {$-$};

    % Outer to inner sum
    \draw[signal] (s1) -- (Co);
    \draw[signal] (Co) -- (s2) node[label, midway]{$T_{in1,SP}$};
    \node[font=\scriptsize] at (s2.north east) {$+$};
    \node[font=\scriptsize] at (s2.south east) {$-$};

    % Inner chain
    \draw[signal] (s2) -- (Ci);
    \draw[signal] (Ci) -- (G1) node[label, midway]{$F_{1,SP}$};
    \draw[signal] (G1) -- (G2) node[label, midway]{$T_{in1}$};

    % Output
    \draw[signal] (G2) -- ++(1.4,0) node[label, pos=1.0]{$T_{out2}$};

    % Inner feedback
    \draw[signal] ($(G1.east)+(0.4,0)$) -- ++(0,-1.5) -| (s2);
    \node[font=\scriptsize, below] at ($(G1.east)+(0.4,-1.5)$) {$T_{in1}$};

    % Outer feedback
    \draw[signal] ($(G2.east)+(0.7,0)$) -- ++(0,-2.5) -| (s1);
    \node[font=\scriptsize, below] at ($(G2.east)+(0.7,-2.5)$) {$T_{out2}$};

    % Brace for inner loop
    \draw[decorate, decoration={brace, amplitude=5pt}, darkblue]
      ($(s2.north west)+(-0.1,0.6)$) -- ($(G1.north east)+(0,0.6)$)
      node[midway, above=6pt, font=\scriptsize, darkblue]{\textbf{Inner loop (fast)}};

    % Brace for outer loop
    \draw[decorate, decoration={brace, amplitude=5pt}, darkblue!50]
      ($(s1.north west)+(-0.1,1.4)$) -- ($(G2.north east)+(0,1.4)$)
      node[midway, above=6pt, font=\scriptsize, darkblue!70]{\textbf{Outer loop (slow)}};
  \end{tikzpicture}
  \caption{Cascade control block diagram: inner PI+P loop ($C_{\text{inner}}$) regulates
    $T_{in1}$; outer P loop ($C_{\text{outer}}$) drives $T_{out2}$ to setpoint.}
  \label{fig:cascade_block}
\end{figure}

\section{Inner Loop --- PI+P Controller}

\subsection{Controller Structure}

The inner loop controller is a discrete-time PI with \textit{proportional action on the
measurement} (P-on-M), also known as PI+P or velocity-form PI. The update equation is:

\begin{equation}
  \Delta u[k] = -K_p\,\bigl(y[k] - y[k-1]\bigr)
                + \frac{K_p}{T_i}\,e[k]\,T_s,
  \label{eq:velocity_increment}
\end{equation}
\begin{equation}
  u[k] = \mathrm{clamp}\bigl(u[k-1] + \Delta u[k],\; F_{1,\min},\; F_{1,\max}\bigr),
  \label{eq:velocity_clamp}
\end{equation}

where $y[k] = T_{in1}$ (measurement), $e[k] = T_{in1,SP} - T_{in1}$ (error), $T_s$
the sampling period (\SI{3}{\second}), and $u[k] = F_{1,SP}$ (controller output in
\si{\metre\cubed\per\hour}).

\subsection{Comparison: Standard PI vs PI+P}

\begin{table}[H]
  \centering
  \caption{Standard PI (position form) vs PI+P (velocity form with P-on-M).}
  \label{tab:pi_comparison}
  \begin{tabular}{@{}L{4.5cm}L{4.5cm}L{4.5cm}@{}}
    \toprule
    \textbf{Property} & \textbf{Standard PI} & \textbf{PI+P (P-on-M)} \\
    \midrule
    Setpoint step response & Proportional kick in $u$ & No kick ($P$ acts on $\Delta y$) \\
    Disturbance rejection  & Good                     & Identical \\
    Integral form          & Explicit sum of errors   & Implicit in increment \\
    Bumpless transfer      & Requires tracking logic  & Natural: initialise $u[k-1]$, $y[k-1]$ \\
    Anti-windup            & Back-calculation needed  & Implicit via output clamping \\
    \bottomrule
  \end{tabular}
\end{table}

The proportional term $-K_p(y[k]-y[k-1])$ acts on the \textit{measurement change} rather
than the error change. For a setpoint step, $y$ does not change instantaneously (only $e$
does), so $\Delta u = 0$ from the proportional term. This prevents the abrupt $F_{1,SP}$
jump that a standard PI would produce, resulting in smoother valve actuation.

\subsection{Anti-Windup}

Output clamping in equation~\eqref{eq:velocity_clamp} provides natural anti-windup
protection. Because $u[k-1]$ is updated to the \textit{clamped} value, the increment
at the next step is computed from the physically actuated output, not from a hypothetical
unclamped value. If $F_{1,SP}$ is saturated at $F_{1,\max}$ and $T_{in1}$ is still below
setpoint ($e[k] > 0$), the integral contribution $\Delta u[k]$ is positive but clamped.
As soon as $T_{in1}$ rises above setpoint ($e[k] < 0$), $\Delta u[k]$ becomes negative and
the output decreases --- without any delay due to windup.

\subsection{Bumpless Transfer}
\label{sec:bumpless}

When the operator switches from MANUAL to AUTO-INNER, the controller must begin at the
current manual $F_{1,SP}$ value with no discontinuity. The \code{init\_output()} method
initialises the velocity form:

\begin{lstlisting}[language=Python,
  caption={Bumpless transfer initialisation in \texttt{cascade\_controller.py}.}]
def init_output(self, current_output: float, current_meas: float):
    """Called before switching from MANUAL to AUTO to prevent bumps."""
    self.last_output = current_output   # sets u[k-1]
    self._last_meas  = current_meas     # sets y[k-1]
\end{lstlisting}

At the first controller cycle after \code{init\_output()}:
\begin{itemize}[itemsep=2pt]
  \item $u[k-1]$ = current manual $F_{1,SP}$ (the value the operator had set)
  \item $\Delta u[k] \approx 0$ from the proportional term (since $y[k] \approx y[k-1]$)
  \item The integral contribution is small if $T_{in1} \approx T_{in1,SP}$
\end{itemize}
The first PI+P output is therefore very close to the manual value --- no bump.

\subsection{Python Implementation}

\begin{lstlisting}[language=Python,
  caption={PI+P velocity form in \texttt{cascade\_controller.py}.}]
class PIplusP:
    def __init__(self, Kp=1.0, Ti=60.0,
                 out_min=0.0, out_max=10.0, dt=3.0):
        self.Kp = Kp
        self.Ti = Ti
        self.out_min = out_min
        self.out_max = out_max
        self.dt = dt
        self.last_output = 0.0
        self._last_meas  = None

    def update(self, setpoint: float, measurement: float) -> float:
        if self._last_meas is None:
            self._last_meas = measurement
        error   = setpoint - measurement
        delta_u = (-self.Kp * (measurement - self._last_meas)
                   + (self.Kp / max(self.Ti, 0.01)) * error * self.dt)
        clamped = max(self.out_min,
                      min(self.out_max, self.last_output + delta_u))
        self.last_output = clamped
        self._last_meas  = measurement
        return clamped
\end{lstlisting}

\section{Outer Loop --- Proportional Controller}

\subsection{Structure}

The outer loop translates the cold outlet setpoint $T_{out2,SP}$ into a hot inlet
setpoint $T_{in1,SP}^{\text{cmd}}$:
\begin{equation}
  T_{in1,SP}^{\text{cmd}}
  = T_{in1,SP}^{\text{base}} + K_{p,ext}\,(T_{out2,SP} - T_{out2,HE1}),
  \label{eq:outer}
\end{equation}

where $T_{in1,SP}^{\text{base}}$ is an operator-defined base setpoint and $K_{p,ext}$
is the outer proportional gain~[\si{\degreeCelsius\per\degreeCelsius}].
Equation~\eqref{eq:outer} is evaluated every 3\,s (outer loop period equal to the
poll interval).

\subsection{Steady-State Offset Analysis}

With a purely proportional outer controller, there is in general a non-zero
steady-state error in $T_{out2}$. At steady state, the inner PI+P loop ensures
$T_{in1} = T_{in1,SP}^{\text{cmd}}$ exactly (due to the integral term).
Substituting into~\eqref{eq:outer} and assuming a linearised outer process gain $K_2$
(i.e.\ $T_{out2,ss} = K_2 T_{in1,ss} + d$):
\begin{equation}
  T_{out2,ss}
  = \frac{K_2\,(T_{in1,SP}^{\text{base}} + K_{p,ext}\,T_{out2,SP}) + d}{1 + K_2 K_{p,ext}}.
  \label{eq:outer_ss}
\end{equation}

The residual error $|T_{out2,SP} - T_{out2,ss}|$ decreases as $K_{p,ext}$ increases,
but an overly large $K_{p,ext}$ risks destabilising the outer loop. In practice, the
operator can eliminate the offset by trimming $T_{in1,SP}^{\text{base}}$.

A pure-integral outer loop (PI outer) would eliminate this offset at the cost of requiring
an outer loop identification experiment. The current design uses P only, with a view to
upgrading to PI once the inner loop has been validated (see Section~\ref{sec:perspectives}).

\section{Operating Modes and State Machine}

The \code{CascadeController} class implements a three-mode state machine. Mode transitions
and their implications are summarised in Table~\ref{tab:modes} and
Figure~\ref{fig:state_machine}.

\begin{table}[H]
  \centering
  \caption{Cascade controller operating modes.}
  \label{tab:modes}
  \begin{tabular}{@{}L{2.8cm}L{1.8cm}L{1.8cm}L{6.5cm}@{}}
    \toprule
    \textbf{Mode} & \textbf{Inner} & \textbf{Outer} & \textbf{Operator input} \\
    \midrule
    MANUAL     & OFF & OFF & $F_{1,SP}$ set directly from the UI \\
    AUTO-INNER & ON  & OFF & $T_{in1,SP}$ set by operator \\
    AUTO-FULL  & ON  & ON  & $T_{out2,SP}$ set by operator \\
    \bottomrule
  \end{tabular}
\end{table}

\begin{figure}[H]
  \centering
  \begin{tikzpicture}[
    state/.style={rounded rectangle, draw=darkblue, thick, fill=blue!8,
                  minimum width=2.8cm, minimum height=1.0cm, font=\small, align=center},
    trans/.style={-Stealth, thick, draw=darkblue},
  ]
    \node[state] (M)  at (0,0)   {MANUAL};
    \node[state] (AI) at (5,0)   {AUTO-INNER};
    \node[state] (AF) at (10,0)  {AUTO-FULL};

    % Forward transitions
    \draw[trans] (M)  to[bend left=20]
      node[above, font=\scriptsize]{init\_output() bumpless} (AI);
    \draw[trans] (AI) to[bend left=20]
      node[above, font=\scriptsize]{outer activates} (AF);

    % Reverse transitions
    \draw[trans] (AI) to[bend left=20]
      node[below, font=\scriptsize]{hold $F_{1,SP}$} (M);
    \draw[trans] (AF) to[bend left=20]
      node[below, font=\scriptsize]{hold $F_{1,SP}$} (AI);
    \draw[trans] (AF) to[out=250, in=290, looseness=1.5]
      node[below, font=\scriptsize]{hold} (M);
  \end{tikzpicture}
  \caption{Cascade controller mode state machine.}
  \label{fig:state_machine}
\end{figure}

\textbf{Identification interlock.} The step test is only permitted when mode~=~MANUAL.
The \code{POST /api/identification/start} endpoint returns an error if mode~$\neq$~MANUAL.
Conversely, mode changes are blocked while a step test is running.

\section{Experimental Validation}

\pendingbox{%
  Section to be completed after closed-loop tests with Michal following the
  identification session of 2026-05-19.%
}

\subsection{Inner Loop --- Setpoint Step Response}

\begin{figure}[H]
  \centering
  \fbox{\begin{minipage}{0.85\textwidth}
    \centering\vspace{2cm}
    \textit{[Figure: $T_{in1,HE1}$ and $F_{1,SP}$ response to a $T_{in1,SP}$ step
    of $+5\,^\circ$C. Compare Normal vs Aggressive IMC tuning.]}
    \vspace{2cm}
  \end{minipage}}
  \caption{Inner loop setpoint step response --- to be completed.}
  \label{fig:inner_step}
\end{figure}

Metrics to report: rise time (10\%--90\%), settling time ($\pm$2\,\%), overshoot (\%),
steady-state error.

\subsection{Inner Loop --- Disturbance Rejection}

\begin{figure}[H]
  \centering
  \fbox{\begin{minipage}{0.85\textwidth}
    \centering\vspace{2cm}
    \textit{[Figure: $T_{in1,HE1}$ and $T_{out2,HE1}$ responses to a $T_{hin}$
    disturbance. MANUAL vs AUTO-INNER comparison.]}
    \vspace{2cm}
  \end{minipage}}
  \caption{Inner loop disturbance rejection --- to be completed.}
  \label{fig:inner_dist}
\end{figure}

\subsection{Full Cascade --- $T_{out2}$ Setpoint Tracking}

\begin{figure}[H]
  \centering
  \fbox{\begin{minipage}{0.85\textwidth}
    \centering\vspace{2cm}
    \textit{[Figure: $T_{out2,HE1}$ and $T_{in1,HE1}$ responses to a $T_{out2,SP}$
    step. Shows outer P controller adjusting $T_{in1,SP}$.]}
    \vspace{2cm}
  \end{minipage}}
  \caption{Full cascade setpoint tracking response --- to be completed.}
  \label{fig:cascade_tracking}
\end{figure}

\subsection{Discussion}

\pendingbox{%
  Quantify steady-state $T_{out2}$ offset (P outer limitation), discuss selected
  tuning level and any stability observations.%
}
